\chapter{Unified Platform Implementation, Azure Deployment, and Business Impact Assessment}

\section*{Introduction}
\addcontentsline{toc}{section}{Introduction}

This chapter documents the FiberGate unified web platform, the third pillar of the FiberGate ecosystem. It brings the SENTINEL automation fleet and the FiberGate AI Service together into a single interface for the Guichet OI: PAR France unit at Orange France, combining an Angular and Spring Boot full-stack application, a FastAPI execution tier, role-based access control, and an Azure cloud deployment.

Section 6.1 covers the architecture and implementation; 6.2 the functional modules; 6.3 security; 6.4 the Azure deployment; 6.5 the CI/CD pipeline; 6.6 the ecosystem integrations; 6.7 testing and validation; 6.8 the business impact assessment; and 6.9 limitations and future work.

\bigskip

\section{Platform Implementation}

The platform is a three-tier web application: an Angular SPA (port 4200), a Spring Boot 3.2.4 orchestration server (port 8080) handling authentication and job routing, and an asynchronous FastAPI execution tier (port 8000) running automation jobs and the AI inference endpoint. All tiers communicate over HTTP with JWT bearer tokens.

\begin{figure}[H]
  \centering
  \fbox{\parbox[c][5cm]{0.55\textwidth}{\centering\itshape Placeholder --- Figure 6.1\\Global system architecture diagram: Angular SPA, Spring Boot orchestration layer, FastAPI execution tier, and Microsoft 365 connections via Power Automate / Graph API}}
  \caption{Global system architecture}
  \label{fig:architecture}
\end{figure}

\subsection{Backend Architecture (Spring Boot)}

The Spring Boot orchestration layer authenticates and authorises every request, proxies automation job submissions to FastAPI, and persists relational data (users, roles, audit entries, job metadata) through JPA over PostgreSQL. It follows a standard layered architecture: controllers expose REST endpoints, services hold business logic, and repositories handle persistence via Spring Data JPA. Cross-cutting concerns (JWT validation, CORS, exception handling, audit logging) are implemented as filters and aspects so they apply uniformly without duplication. As the security boundary, Spring Boot validates every JWT and resolves the user role before forwarding verified requests to FastAPI, which therefore only ever receives authenticated traffic.

\subsection{Frontend Architecture (Angular)}

The frontend is an Angular 17 SPA with a component-driven architecture: each functional module (authentication, operator dashboard, automation panels, AI review interface, administration console) is an independently routed module with its own components, services, and route guards. State is managed through RxJS observables so dashboard panels react to job-status changes without polling. An HTTP interceptor attaches the JWT to every request, and route guards enforce role-based navigation before any server call is made. The interface was designed from high-fidelity Figma prototypes validated with stakeholders beforehand, and adapts to the authenticated role: Consultants see their dossier queue and AI review tasks, Pilots see automation controls and the distribution upload, and the SDM sees the audit trail and performance dashboard.

\begin{figure}[H]
  \centering
  \fbox{\parbox[c][5cm]{0.5\textwidth}{\centering\itshape Placeholder --- Screenshot\\Angular operator dashboard (role-adapted view)}}
  \caption{FiberGate operator dashboard interface}
  \label{fig:dashboard-screenshot}
\end{figure}

\subsection{API Design and Service Integration}

The REST API is organised by resource: identity/session under \texttt{/api/auth/}, business resources (jobs, dossiers, users) under \texttt{/api/jobs/}, \texttt{/api/dossiers/}, \texttt{/api/users/}, and administration under \texttt{/api/admin/} and \texttt{/api/audit/}, with conventional HTTP verbs throughout.

The service integration layer decouples Angular from the execution backends through an asynchronous job pattern: Spring Boot validates a submission, records it, and forwards it to FastAPI, which starts a background task, returns a UUID immediately, and lets the dashboard poll \texttt{/api/jobs/\{id\}/status} for progress. This non-blocking design keeps long-running jobs (e.g.\ a 112-task distribution batch taking up to 3 min 40 s) from blocking the UI or timing out the load balancer.

\begin{figure}[H]
  \centering
  \fbox{\parbox[c][4.5cm]{0.55\textwidth}{\centering\itshape Placeholder --- Figure 6.2\\Sequence diagram: asynchronous job submission flow (Angular $\rightarrow$ Spring Boot $\rightarrow$ FastAPI $\rightarrow$ polling)}}
  \caption{Asynchronous job submission and polling sequence}
  \label{fig:job-sequence}
\end{figure}

\subsection{Database Architecture and Schema}

The platform uses a PostgreSQL database under Liquibase migrations, organised into five entity groups: \emph{identity} (\texttt{users}, \texttt{roles}, \texttt{user\_roles}, linked to Microsoft Entra ID); \emph{jobs} (\texttt{automation\_jobs}, \texttt{job\_logs}, recording every execution and its status); \emph{dossiers} (\texttt{dossiers}, \texttt{document\_extractions}, \texttt{field\_corrections}, holding AI results and consultant corrections that feed the Chapter 5 retraining loop); \emph{audit} (\texttt{audit\_log}, recording every state-changing operation with before/after values, feeding the SDM audit view); and \emph{configuration} (\texttt{system\_config}, runtime-editable parameters such as confidence thresholds and model version).

\begin{figure}[H]
  \centering
  \fbox{\parbox[c][5.5cm]{0.6\textwidth}{\centering\itshape Placeholder --- Figure 6.3\\Entity-relationship diagram of the five schema groups (identity, jobs, dossiers, audit, configuration)}}
  \caption{Relational database schema overview}
  \label{fig:db-schema}
\end{figure}

\section{Functional Modules and Platform Features}

The platform provides five role-guarded functional modules sharing a common Angular shell (navigation bar, notification centre, session indicator).

\subsection{Authentication and Session Management}

The module implements a stateless JWT session model: on login, Spring Boot verifies credentials, issues a signed JWT with role claims, and the Angular client stores it in memory (not \texttt{localStorage}, to mitigate XSS token theft). The interceptor attaches the token to every request, a silent refresh keeps sessions alive, and logout invalidates the refresh record server-side. The login page is the only unauthenticated route; all others are protected by route guards that redirect on missing or insufficient claims.

\begin{figure}[H]
  \centering
  \fbox{\parbox[c][4.5cm]{0.45\textwidth}{\centering\itshape Placeholder --- Screenshot\\Login / authentication screen}}
  \caption{Authentication screen}
  \label{fig:login-screenshot}
\end{figure}

\subsection{User and Role Administration}

Available only to the SDM role, this module provides a searchable user table with inline role assignment, account activation/deactivation, and password resets via Exchange. New users are provisioned by Microsoft Entra ID identifier, with the Graph API resolving the display name and email automatically. Role enforcement is layered: the Angular guard blocks navigation, the Spring Boot filter rejects insufficiently privileged API calls, and the database permission model restricts data access even if the upper layers fail.

\begin{figure}[H]
  \centering
  \fbox{\parbox[c][4.5cm]{0.5\textwidth}{\centering\itshape Placeholder --- Screenshot\\User and role administration console}}
  \caption{Administration console}
  \label{fig:admin-screenshot}
\end{figure}

\subsection{Automation Monitoring and Control Dashboard}

The main interface for Pilots, this dashboard presents one panel per SENTINEL agent showing job status, a submission form, and a live log view streamed from the FastAPI job store. The POI panel (Agent 1) reports scanned emails, extracted fields, and created SUIVI rows / Planner tasks; the CAFF panel (Agent 2) shows the latest processed return and its routing outcome; the distribution panel (Agent 3) accepts the daily Excel batch, shows progress by demand type, and links to the colour-coded output workbook. A summary bar aggregates the day's job counts, success rates, and processing time across all three agents.

\begin{figure}[H]
  \centering
  \fbox{\parbox[c][5cm]{0.55\textwidth}{\centering\itshape Placeholder --- Screenshot\\Automation monitoring dashboard with per-agent panels and live log view}}
  \caption{Automation monitoring and control dashboard}
  \label{fig:automation-dashboard-screenshot}
\end{figure}

\subsection{AI Service Interaction and Results Management}

The primary Consultant workspace: when a dossier is analysed, the FiberGate AI Service classifies each document (87.7\% accuracy) and extracts the twelve target fields at the reliability reported in Table 5.10. Results appear in a structured review form with, per field, the extracted value, a colour-coded confidence indicator (green/amber/red), and an editable correction box. A completeness verdict (\textsc{demande complete} / \textsc{incomplete}) appears at the top with the triggering fields; the consultant can confirm, override with justification, or request reanalysis. Every correction is stored with its original OCR tokens and coordinates, feeding the continuous-improvement loop, and a download button produces the pre-filled 27-column Excel artefact for Agilis submission.

\begin{figure}[H]
  \centering
  \fbox{\parbox[c][5.5cm]{0.6\textwidth}{\centering\itshape Placeholder --- Screenshot\\AI results review interface: extracted fields, confidence indicators, completeness verdict}}
  \caption{AI results review interface}
  \label{fig:ai-review-screenshot}
\end{figure}

\subsection{System Administration and Audit Trail}

The audit module gives the SDM a chronological, filterable log of every state-changing event (logins, role changes, job submissions, AI confirmations, corrections, configuration edits), each entry showing actor, action, resource, timestamp, and before/after values. A system health panel reports the active model version (read from the AI Service \texttt{/health} endpoint), confidence thresholds, database pool status, and tier availability; thresholds and rate-limit delays can be edited inline, with each change written to \texttt{system\_config} and logged.

\begin{figure}[H]
  \centering
  \fbox{\parbox[c][4.5cm]{0.5\textwidth}{\centering\itshape Placeholder --- Screenshot\\Audit trail and system health panel}}
  \caption{Audit trail and system administration panel}
  \label{fig:audit-screenshot}
\end{figure}

\section{Security Architecture and Access Management}

Security is enforced at every layer of the request lifecycle through three mechanisms: JWT authentication, role-based access control, and a GDPR-aligned data protection strategy.

\subsection{Authentication (JWT)}

Each token is signed with HMAC-SHA256 (256-bit secret) and carries the subject, role list, issue time, and expiry (15 minutes for access tokens, 7 days for refresh tokens), transmitted exclusively over HTTPS. The Spring Boot security filter chain validates the signature and expiry, resolves role claims, and builds the authentication context before any controller runs; invalid or missing tokens receive an immediate HTTP 401.

\subsection{Role-Based Access Control}

Three roles map directly to the unit's operational actors: \textbf{Consultant} (AI review interface, dossier queue, own correction history), \textbf{Pilot} (Consultant access plus automation dashboard and distribution upload, with aggregate but not individual statistics), and \textbf{SDM} (full access: provisioning, role assignment, audit trail, configuration). Enforcement is layered across Angular route guards, Spring Boot \texttt{@PreAuthorize} annotations, and Hibernate row-level filters, so a frontend misconfiguration alone cannot expose protected data.

\begin{figure}[H]
  \centering
  \fbox{\parbox[c][4.5cm]{0.55\textwidth}{\centering\itshape Placeholder --- Figure 6.4\\Role-based access control matrix: Consultant / Pilot / SDM permissions across modules and enforcement layers}}
  \caption{RBAC permission matrix}
  \label{fig:rbac-matrix}
\end{figure}

\subsection{Data Protection and GDPR Compliance}

The platform handles personal data (names, emails, phone numbers) and applies four compliance principles: \emph{data minimisation} (extracted values are stored, but original document images are processed in memory and discarded — never written to the application database); \emph{on-premises annotation} (the Chapter 5 Label Studio server ran on-premises behind a local CORS proxy, so no personal data left the organisation during corpus construction); \emph{audit logging} (every access to personal data is recorded with actor and timestamp, supporting GDPR Article 15 access requests); and a \emph{retention policy} (correction records used for retraining are pseudonymised before storage).

\bigskip

\section{Azure Cloud Deployment}

The deployment target evolved from the VM/App-Service design of Section~3.4 to a fully containerised topology, chosen for three reasons: Docker images eliminate environment drift across slots; Azure Container Apps offers scale-to-zero and per-revision blue-green rollback unavailable in the shared subscription's App Service plans; and a managed PostgreSQL Flexible Server matches the open-source stack while staying behind a private endpoint. The as-built infrastructure below supersedes the Section~3.4 target design.

\begin{figure}[H]
  \centering
  \fbox{\parbox[c][5.5cm]{0.6\textwidth}{\centering\itshape Placeholder --- Figure 6.5\\Azure deployment topology: Container Apps environments, VNet, managed PostgreSQL, Blob Storage, and Microsoft 365 connectivity}}
  \caption{Azure cloud deployment architecture}
  \label{fig:azure-architecture}
\end{figure}

\subsection{Infrastructure Provisioning}

Each tier runs as a Docker container on Azure Container Apps: Nginx serves the Angular SPA, Spring Boot runs on a standard container (2 vCPU / 4 GB RAM), and the GPU-enabled FastAPI/LayoutLMv3 tier requires roughly 8 GB VRAM for the 7 GB model. Images are stored in Azure Container Registry tagged with the Git commit SHA for traceability. PostgreSQL runs on Azure Database for PostgreSQL Flexible Server with automated backups, point-in-time restore, and a private endpoint that blocks public access. Model checkpoints and generated Excel artefacts are stored in versioned, dated folders in Azure Blob Storage, with the active checkpoint selected at startup via the \texttt{system\_config} version identifier.

\subsection{Network Architecture}

All Container Apps run inside a shared VNet and communicate over private DNS, never crossing the public internet; only the Angular CDN origin (Azure Front Door) and the Spring Boot API (Application Gateway with WAF) are publicly exposed. The on-premises Exchange Server is reached by Agents 1 and 2 through a hybrid VPN/ExpressRoute circuit without exposing it to the internet. Microsoft 365 services (SharePoint, Planner, Teams, Entra ID) are accessed via the Graph API using managed identities backed by Key Vault, so no secrets appear in configuration or environment variables.

\subsection{Scalability and High Availability}

The Angular and Spring Boot containers scale to zero off-hours and add replicas above 70\% CPU or a queue of more than ten pending jobs. The GPU inference container keeps one replica always active (a 7 GB model load is too slow for an interactive workflow on cold start), with a second available on demand. Zone redundancy covers the Container Apps environments, and the managed PostgreSQL instance fails over to a standby replica in under two minutes.

\bigskip

\section{CI/CD Pipeline and DevOps Strategy}

Source control, build automation, and release management are split across GitHub and Azure DevOps.

\begin{figure}[H]
  \centering
  \fbox{\parbox[c][5cm]{0.6\textwidth}{\centering\itshape Placeholder --- Figure 6.6\\CI/CD pipeline diagram: GitHub Actions build/test stages $\rightarrow$ Azure DevOps release pipeline $\rightarrow$ Dev / Staging / Production promotion}}
  \caption{CI/CD pipeline and environment promotion flow}
  \label{fig:cicd-pipeline}
\end{figure}

\subsection{Source Control and Branching (GitHub)}

A single monorepo holds the Angular, Spring Boot, FastAPI, and infrastructure-as-code code, under a simplified GitFlow: \texttt{main} is production-protected (PR-only), \texttt{develop} is the integration branch, feature branches follow \texttt{feature/\{ticket-id\}-\{description\}}, and \texttt{hotfix/\{id\}} branches are cut directly from \texttt{main}. Merges to \texttt{develop} or \texttt{main} require an approving review (rotated via CODEOWNERS) and a passing CI build.

\subsection{Continuous Integration (GitHub Actions)}

Every push and pull request triggers dependency installation, static analysis (ESLint, SpotBugs, Ruff), unit tests with coverage, and a Docker build — any failing step blocks the merge. Coverage gates require at least 70\% for the Spring Boot service layer and 60\% for FastAPI route handlers, with reports linked in the PR summary. Images are tagged with branch name and commit SHA; feature-branch images go to the ACR development namespace only.

\subsection{Continuous Delivery (Azure DevOps)}

A release pipeline triggers automatically on image pushes to the staging namespace (merge to \texttt{develop}) or production namespace (merge to \texttt{main}). It applies Terraform plans, updates the ACA revision, runs smoke tests, and notifies the team via Microsoft Teams. Terraform state lives in Azure Storage with locking, and infrastructure changes require their own reviewed PR with a \texttt{terraform plan} step.

\subsection{Environment Promotion}

Three independent environments — each with its own database, storage, and managed identity, configured via Key Vault references so the same image runs unmodified everywhere — form the promotion chain. \textbf{Development} receives every \texttt{develop} merge, runs on synthetic data, and scales to zero off-hours. \textbf{Staging} receives release candidates at production-equivalent sizing for performance testing and user acceptance testing. \textbf{Production} is promoted after UAT sign-off via a manual SDM-gated approval, using a blue-green rollout that ramps a new revision from 10\% to full traffic over five minutes of monitored operation.

\bigskip

\section{Ecosystem Integration Architecture}

The platform connects the three SENTINEL agents, the FiberGate AI Service, and the Microsoft 365 tools used daily by the Guichet OI team.

\begin{figure}[H]
  \centering
  \fbox{\parbox[c][5.5cm]{0.6\textwidth}{\centering\itshape Placeholder --- Figure 6.7\\Ecosystem integration diagram: SENTINEL agents, AI Service, Spring Boot/FastAPI proxy layer, and Microsoft 365 (SharePoint, Planner, Exchange, Graph API)}}
  \caption{Ecosystem integration architecture}
  \label{fig:ecosystem-integration}
\end{figure}

\subsection{SENTINEL Fleet Integration}

The three agents run as Python modules inside the FastAPI tier. Spring Boot proxies job submissions to FastAPI, which runs them as background tasks, stores them in a UUID job store, and exposes a status endpoint for polling. The Angular client never talks to FastAPI directly — every request is authenticated and authorised by Spring Boot first, so FastAPI operates entirely within the trusted network without its own auth layer. Captured stdout/stderr from each run is exposed through the job store, powering the dashboard's live log view.

\subsection{FiberGate AI Service Integration}

When a dossier is submitted, Spring Boot stores the files in Blob Storage and notifies FastAPI, which runs Tesseract 5 OCR, feeds tokens and layout coordinates to LayoutLMv3 in a single forward pass, applies the confidence-gated policy, and returns structured JSON. Spring Boot persists the results to \texttt{document\_extractions} and notifies the frontend; consultant corrections are written to \texttt{field\_corrections} for the next retraining cycle. The active model version, reported via \texttt{/health}, is visible to the SDM in the administration module.

\subsection{Microsoft 365 Integration}

SharePoint SUIVI is the authoritative dossier list, Planner holds task assignments, and Exchange delivers the Agilis/GedAffaire trigger emails. Two mechanisms bridge the platform to this ecosystem: \emph{Power Automate cloud flows} handle high-frequency, event-driven updates — for example, Agent 1's webhook triggers a seven-step flow that updates SUIVI, creates a Planner task, and assigns it via an Entra ID lookup, while Agent 2 uses a simpler dedicated flow for CAFF routing; and \emph{Microsoft Graph API} calls, authenticated via managed identities with minimal OAuth2 scopes (\texttt{Mail.Send}, \texttt{Tasks.ReadWrite}, \texttt{Sites.ReadWrite.All}), handle lower-frequency operations such as directory lookups, Teams notifications, and SharePoint queries.

\subsection{End-to-End Workflow}

The integrations combine into a complete LOC PAR lifecycle: Agent 1 detects the Agilis confirmation email, extracts six fields, and triggers the SUIVI/Planner update; in parallel, Agent 3 classifies the daily 112-file batch, assigns consultants by continuity-preserving round robin, and dispatches Planner tasks via Power Automate in under four minutes. Agent 2 routes each CAFF return within sixty seconds regardless of arrival time, eliminating the previous 1-day KPI miss. For completeness validation, the consultant uploads documents to the AI interface, which classifies and extracts fields in under thirty seconds per dossier, reducing what was a manual reading and entry task to a review-and-confirm interaction.

\begin{figure}[H]
  \centering
  \fbox{\parbox[c][5cm]{0.6\textwidth}{\centering\itshape Placeholder --- Figure 6.8\\End-to-end workflow diagram: from Agilis/GedAffaire email or PAR batch to Planner task and Agilis-ready Excel artefact}}
  \caption{End-to-end integrated business workflow}
  \label{fig:e2e-workflow}
\end{figure}

\section{System Testing and Validation}

The platform was validated at three levels: unit/integration testing, user acceptance testing, and cloud deployment validation.

\subsection{Unit and Integration Testing}

Spring Boot unit tests (JUnit 5, Mockito) reach 74\% line coverage across 142 cases; FastAPI agent tests (Pytest) reach 63\% across 89 cases, with most gaps in exception-handling paths for malformed Exchange responses. Integration tests use an embedded H2 database with MockMvc for Spring Boot and a mock Exchange Server for the agents; three end-to-end tests cover the job-submission round trip, the document-analysis pipeline, and Power Automate webhook delivery to SharePoint.

\subsection{User Acceptance Testing}

Five participants (two Consultants, two Pilots, the SDM) ran structured test scripts across two three-hour sessions on Staging. Findings: confidence colour-coding was understood without explanation by all participants; live log streaming was rated highly useful by Pilots; a usability gap in the distribution upload (expected drag-and-drop, found only a file browser) was fixed by adding a drop zone before production; and the SDM's request for a resource-type filter in the audit trail was added as a dropdown before sign-off. All critical/high-priority issues were resolved before the production gate.

\subsection{Cloud Deployment Validation}

Infrastructure provisioning was validated by a clean Terraform apply and a 34-point manual checklist (ACA environments, container revisions, PostgreSQL, Blob Storage, Key Vault, VNet peerings). Under a simulated load of twenty concurrent users, Spring Boot held a 210 ms median response time, the GPU inference endpoint processed one dossier every 28 seconds (up to two concurrent without queuing), and the 112-task distribution batch completed in 3 min 40 s — within the five-minute requirement. Failover tests confirmed a sub-90-second PostgreSQL standby promotion with no data loss, and a deliberately broken image triggered an automatic blue-green rollback within three minutes.

\bigskip

\section{Business Impact Assessment and Return on Investment}

\subsection{Operational Improvement Metrics}

Table 6.1 summarises measured performance before and after deployment.

\begin{table}[H]
  \centering
  \caption{Operational improvement metrics by automated workflow.}
  \renewcommand{\arraystretch}{1.0}
  \footnotesize
  \begin{tabularx}{\textwidth}{|X|X|X|X|X|}
    \hline
    \textbf{System} & \textbf{Trigger} & \textbf{Technology} & \textbf{Before / After} & \textbf{Reduction} \\ \hline
    Agent 1: POI Validation & Agilis confirmation email & Python, exchangelib, Power Automate (7 steps) & 15 to 20 min / under 60 s & approximately 95\% \\ \hline
    Agent 2: CAFF Returns & GedAffaire mail & Python, exchangelib, atomic Power Automate flow & KPI missed / guaranteed & 100\% of misses eliminated \\ \hline
    Agent 3: Distribution & Daily Excel batch of 112 files & pandas, openpyxl, round-robin with Planner lookup & 30 to 45 min / under 2 min & approximately 95\% \\ \hline
    FiberGate AI & PDF or PNG dossier & LayoutLMv3, Tesseract 5 OCR & Manual review / 87.7\% class.\ accuracy, 0.448 span-F1 & review-assisted \\ \hline
  \end{tabularx}
\end{table}

The largest gain is in POI validation: a sequence that took a consultant 15--20 minutes (read, extract six fields, update SUIVI, create and assign the Planner task) now completes in under sixty seconds, saving roughly five analytical hours per day at the unit's volume of 112 files. CAFF routing, previously prone to being missed after 5 PM and causing KPI breaches, is now guaranteed within sixty seconds of arrival regardless of time of day.

\begin{figure}[H]
  \centering
  \fbox{\parbox[c][4.5cm]{0.55\textwidth}{\centering\itshape Placeholder --- Figure 6.9\\Bar chart: per-task processing time before vs.\ after automation, by workflow}}
  \caption{Before/after processing-time comparison}
  \label{fig:impact-chart}
\end{figure}

\subsection{Time and Cost Savings}

The fifteen-consultant unit currently spends an average of forty-five minutes per file (roughly twenty minutes on POI processing, ten on CAFF routing, and sixteen to twenty-four seconds per file on distribution). Automation reduces each of these to a two-minute review-and-confirm step. At a fully loaded cost of 45 €/hour and 112 files/day across 220 working days, the projected annual saving is approximately \textbf{12,320 hours} of analytical work — about \textbf{554,400 €} per year, or roughly 37,000 € per consultant, with the development cost recovered within the first two months of production.

\subsection{Workforce Productivity Gains}

Before deployment, about 55\% of consultant time went to repetitive low-complexity tasks (reading emails, copying fields, creating tasks, distributing work); that time is now redirected to analytical and relational work that cannot be automated. The AI review interface cuts per-dossier validation time from an estimated fifteen minutes to three to five minutes for well-structured (Strong-tier) cases. UAT participants confirmed feeling more in control of their workload, and Pilots specifically valued the new operational visibility from the live dashboard.

\subsection{Strategic Value and Transferability}

The FiberGate blueprint is reusable beyond this project through four components: a domain-agnostic three-tier platform (Angular/Spring Boot/FastAPI) that can host any agent-and-AI combination behind a unified interface; a lightweight webhook-to-cloud-flow integration pattern usable by any Microsoft 365–based organisation; an adaptable LayoutLMv3 fine-tuning pipeline for other structured-document corpora; and a confidence-gated, correction-capture deployment policy suited to regulated environments requiring human review. Orange France has expressed interest in extending the model to other DINFRA units, and the Amaris Consulting team has identified three additional Orange departments and two external clients as candidates for moderate-effort adaptation.

\bigskip

\section{System Limitations and Future Research Perspectives}

\textbf{Tabular field performance.} Three field types (\texttt{Nb\_log\_res}, \texttt{Nb\_log\_pro}, \texttt{Nombre\_Logement\_Lot\_MacroLot}) score zero span-F1 (Table 5.10), because LayoutLMv3's 512-token window cannot cover both a table's header row and its target cell in dense layouts. A longer-context model or a table-region pre-processing step is needed.

\textbf{OCR dependency.} Extraction quality is bounded by Tesseract 5 output; heavily stamped, low-contrast, or skewed pages degrade tokens and propagate errors. A pre-processing stage (deskew, denoise, contrast correction) and evaluation of a commercial OCR engine for low-quality documents are recommended.

\textbf{Single-page processing.} Multi-page documents cannot be analysed in one inference call, and there is no cross-page attention. A hierarchical document representation or an overlapping sliding window would address this.

\textbf{Exchange Server dependency.} Agents 1 and 2 rely on EWS/NTLM to an on-premises Exchange Server, requiring manual reconfiguration on certificate or mailbox changes. Migrating to Exchange Online with Graph API access would remove the VPN dependency and align this integration with the rest of the Azure-native architecture.

\textbf{Retraining loop maturity.} The correction-capture and retraining framework (Chapter 5, Section 6.6.2) is fully implemented but has not yet completed a production cycle; the first cycle, expected after roughly two months of correction accumulation, will provide the first evidence of whether it can lift Critical-tier fields above the operational threshold.

Future directions include few-shot prompting with large multimodal models as an annotation-free alternative to fine-tuning, active learning to prioritise which corrections to annotate first, and an explainability layer that highlights the document regions driving each prediction.

\section*{Conclusion}
\addcontentsline{toc}{section}{Conclusion}

This chapter documented the FiberGate unified web platform and its connections to the SENTINEL fleet, the AI Service, and Microsoft 365. The three-tier Angular/Spring Boot/FastAPI architecture cleanly separates presentation, orchestration, and execution; the JWT-based RBAC model keeps each actor — Consultant, Pilot, SDM — within a well-defined permission boundary; and the Azure deployment, driven by a GitHub Actions / Azure DevOps CI/CD pipeline with blue-green releases and infrastructure as code, delivers a reliable production service integrated into the team's existing SharePoint and Planner habits.

The measured business impact is substantial: a 95\% reduction in POI validation and distribution processing time, full elimination of the 1-day CAFF KPI miss, and a projected annual saving of approximately 554,400 € with the investment recovered within two months of going live. Combining automation agents, document AI, and a unified web platform over Microsoft 365, the FiberGate model offers a replicable blueprint for administrative automation in telecoms — with concrete extension opportunities already identified at three Orange France departments and two external clients.
